\begin{abstract}
The software around a language model in an agent is called a scaffold, a framework, a runtime or a harness, and none of these names says what it contains. We take the position that an LLM agent is a model wrapped in a harness, $\textit{Agent} = \textit{Harness}(\textit{Model})$, and that every harness is built from five mechanistic primitives: \prim{Input}, \prim{Context}, \prim{Model interface}, \prim{Output} and \prim{Control flow}. Each is defined by what it does and by what it never does, so the five are mutually exclusive. Everything else said about a harness is at one of two other levels: a \emph{design decision} within a primitive, which is where harnesses differ, or a \emph{production concern} (persistence, safety and security, observability, resilience, performance, evaluation), a question asked of every primitive rather than a part of the harness. As evidence, we deconstruct six harnesses from their primary sources at pinned versions: Claude Code, Codex CLI, opencode, pi, Hermes Agent and a LangGraph workflow. Every component found has a home in one primitive or one concern, and the differences between the six are design decisions. The vocabulary gives a procedure for locating a failure in one of five places before attributing it to the model, and a schema for reporting the harness alongside the model. The evidence for every assignment, and the minimal harnesses the definitions were drawn from, are published with the paper.
\end{abstract}
\begin{IEEEkeywords}
LLM agents, agent harness, scaffold, agent architecture, software engineering agents
\end{IEEEkeywords}
